% GENERATED FILE. Edit canonical lessons, then run npm run generate:books.
% canonical-source-hash: fnv1a64:0735f383fcc24c80
% canonical-lessons: GU-C327-utsav, GU-C327-ujavvu, GU-C327-parampara, GU-C327-rivaj, GU-C327-parva

\chapter{The Story Behind the Festival}
\label{ch:gu-the-story-behind-the-festival}
\hlchaptermodality{327}

\begin{chapteropening}
\textbf{By the end of this chapter:} \emph{I can say how a festival is kept and why.}

Complete the last lesson of chapter 327: I can say how a festival is kept and why.
\end{chapteropening}

\section[\texorpdfstring{\gu{ઉત્સવ} (utsav)}{ઉત્સવ (utsav)}]{\gu{ઉત્સવ} (utsav) — a celebration}

\label{lesson:GU-C327-utsav}



\begin{quote}
Before the new one: say the Gujarati for a border, then the Gujarati for to found.
\end{quote}

\subsection*{You'll want to know: \gu{ઉત્સવ}}
\textbf{\gu{ઉત્સવ}} — \emph{utsav} — \textquotedblleft{}a celebration\textquotedblright{}.

\begin{quote}
\gu{ગામમાં} \gu{ત્રણ} \gu{દિવસનો} \gu{ઉત્સવ} \gu{હતો}.
\end{quote}

\begin{quote}
\emph{gāmmā̃ traṇ divasno utsav hato.}
\end{quote}

\begin{quote}
There was a three-day celebration in the village.
\end{quote}

A masculine noun.

\begin{cousinweb}
is a tatsama, Sanskrit \emph{utsava}, \textquotedblleft{}a festival, a joyful occasion\textquotedblright{} (Monier-Williams). \textbf{\gu{તહેવાર}} is a festival of the calendar; a \textbf{\gu{ઉત્સવ}} is any celebration, big or small.
\end{cousinweb}

\begin{grammarlens}[title={What you've built}]
One more word for telling what happened.
\end{grammarlens}

\subsection*{Your turn}
\begin{itemize}
\raggedright
  \item \emph{Say it:} \emph{utsav}
  \item \emph{Say it:} \emph{utsav}, once more
  \item \emph{Say it:} \emph{gāmmā̃ traṇ divasno utsav hato.} — There was a three-day celebration in the village.
  \item \emph{Recall:} say the Gujarati for a border, then the Gujarati for to found, then say \emph{utsav} again
\end{itemize}

\begin{tcolorbox}[breakable,colback=teal!4,colframe=teal!35!black,title={Before you move on}]
What does \gu{ઉત્સવ} mean? (\textquotedblleft{}A celebration\textquotedblright{}.) Say it once more.
\end{tcolorbox}
\section[\texorpdfstring{\gu{ઉજવવું} (ujavvũ)}{ઉજવવું (ujavvũ)}]{\gu{ઉજવવું} (ujavvũ) — to celebrate}

\label{lesson:GU-C327-ujavvu}



\begin{quote}
Before the new one: say the Gujarati for to found, then the Gujarati for a celebration.
\end{quote}

\subsection*{You'll want to know: \gu{ઉજવવું}}
\textbf{\gu{ઉજવવું}} — \emph{ujavvũ} — \textquotedblleft{}to celebrate\textquotedblright{}.

\begin{quote}
\gu{અમે} \gu{દિવાળી} \gu{ઘરે} \gu{ઉજવી}.
\end{quote}

\begin{quote}
\emph{ame divāḷī ghare ujvī.}
\end{quote}

\begin{quote}
We celebrated Diwali at home.
\end{quote}

Past, with \gu{તેણે} (he or she …): \gu{તેણે} \gu{તહેવાર} \gu{ઉજવ્યો}, \emph{teṇe tahevār ujavyo}.

\begin{cousinweb}
\textbf{\gu{ઉજવવું}} (\emph{ujavvũ}) comes from \emph{*udyāpayati}, a form CDIAL reconstructs, where it lists the Gujarati verb, \textquotedblleft{}to conclude a religious vow\textquotedblright{} (CDIAL 2053): the feast at the end of a vow became any celebration. It agrees with the festival: \textbf{\gu{દિવાળી} \gu{ઉજવી}}.
\end{cousinweb}

\begin{grammarlens}[title={What you've built}]
One more word for telling what happened.
\end{grammarlens}

\subsection*{Your turn}
\begin{itemize}
\raggedright
  \item \emph{Say it:} \emph{ujavvũ}
  \item \emph{Say it:} \emph{ujavvũ}, once more
  \item \emph{Say it:} \emph{ame divāḷī ghare ujvī.} — We celebrated Diwali at home.
  \item \emph{Recall:} say the Gujarati for to found, then the Gujarati for a celebration, then say \emph{ujavvũ} again
\end{itemize}

\begin{tcolorbox}[breakable,colback=teal!4,colframe=teal!35!black,title={Before you move on}]
What does \gu{ઉજવવું} mean? (\textquotedblleft{}To celebrate\textquotedblright{}.) Say it once more.
\end{tcolorbox}
\section[\texorpdfstring{\gu{પરંપરા} (paramparā)}{પરંપરા (paramparā)}]{\gu{પરંપરા} (paramparā) — a tradition}

\label{lesson:GU-C327-parampara}



\begin{quote}
Before the new one: say the Gujarati for a celebration, then the Gujarati for to celebrate.
\end{quote}

\subsection*{You'll want to know: \gu{પરંપરા}}
\textbf{\gu{પરંપરા}} — \emph{paramparā} — \textquotedblleft{}a tradition\textquotedblright{}.

\begin{quote}
\gu{આ} \gu{અમારા} \gu{ગામની} \gu{જૂની} \gu{પરંપરા} \gu{છે}.
\end{quote}

\begin{quote}
\emph{ā amārā gāmnī jūnī paramparā che.}
\end{quote}

\begin{quote}
This is an old tradition of our village.
\end{quote}

A feminine noun.

\begin{cousinweb}
\textbf{\gu{પરંપરા}} (\emph{paramparā}) is a tatsama, Sanskrit \emph{paramparā}, \textquotedblleft{}an unbroken series, a succession\textquotedblright{} (Monier-Williams): what is handed down from one generation to the next.
\end{cousinweb}

\begin{grammarlens}[title={What you've built}]
One more word for telling what happened.
\end{grammarlens}

\subsection*{Your turn}
\begin{itemize}
\raggedright
  \item \emph{Say it:} \emph{paramparā}
  \item \emph{Say it:} \emph{paramparā}, once more
  \item \emph{Say it:} \emph{ā amārā gāmnī jūnī paramparā che.} — This is an old tradition of our village.
  \item \emph{Recall:} say the Gujarati for a celebration, then the Gujarati for to celebrate, then say \emph{paramparā} again
\end{itemize}

\begin{tcolorbox}[breakable,colback=teal!4,colframe=teal!35!black,title={Before you move on}]
What does \gu{પરંપરા} mean? (\textquotedblleft{}A tradition\textquotedblright{}.) Say it once more.
\end{tcolorbox}
\section[\texorpdfstring{\gu{રિવાજ} (rivāj)}{રિવાજ (rivāj)}]{\gu{રિવાજ} (rivāj) — a custom}

\label{lesson:GU-C327-rivaj}



\begin{quote}
Before the new one: say the Gujarati for to celebrate, then the Gujarati for a tradition.
\end{quote}

\subsection*{You'll want to know: \gu{રિવાજ}}
\textbf{\gu{રિવાજ}} — \emph{rivāj} — \textquotedblleft{}a custom\textquotedblright{}.

\begin{quote}
\gu{લગ્નનો} \gu{આ} \gu{રિવાજ} \gu{ખૂબ} \gu{જૂનો} \gu{છે}.
\end{quote}

\begin{quote}
\emph{lagnano ā rivāj khūb jūno che.}
\end{quote}

\begin{quote}
This wedding custom is very old.
\end{quote}

A masculine noun.

\begin{cousinweb}
\textbf{\gu{રિવાજ}} (\emph{rivāj}) is usually traced to Persian \emph{rivāj}, \textquotedblleft{}currency, custom\textquotedblright{}, a word Persian took from Arabic: what is current among people. It is masculine: \textbf{\gu{જૂનો} \gu{રિવાજ}}.
\end{cousinweb}

\begin{grammarlens}[title={What you've built}]
One more word for telling what happened.
\end{grammarlens}

\subsection*{Your turn}
\begin{itemize}
\raggedright
  \item \emph{Say it:} \emph{rivāj}
  \item \emph{Say it:} \emph{rivāj}, once more
  \item \emph{Say it:} \emph{lagnano ā rivāj khūb jūno che.} — This wedding custom is very old.
  \item \emph{Recall:} say the Gujarati for to celebrate, then the Gujarati for a tradition, then say \emph{rivāj} again
\end{itemize}

\begin{tcolorbox}[breakable,colback=teal!4,colframe=teal!35!black,title={Before you move on}]
What does \gu{રિવાજ} mean? (\textquotedblleft{}A custom\textquotedblright{}.) Say it once more.
\end{tcolorbox}
\section[\texorpdfstring{\gu{પર્વ} (parva)}{પર્વ (parva)}]{\gu{પર્વ} (parva) — a holy day; a festival day}

\label{lesson:GU-C327-parva}



\begin{quote}
Before the new one: say the Gujarati for a tradition, then the Gujarati for a custom.
\end{quote}

\subsection*{You'll want to know: \gu{પર્વ}}
\textbf{\gu{પર્વ}} — \emph{parva} — \textquotedblleft{}a holy day; a festival day\textquotedblright{}.

\begin{quote}
\gu{દિવાળી} \gu{પ્રકાશનું} \gu{પર્વ} \gu{છે}.
\end{quote}

\begin{quote}
\emph{divāḷī prakāśnũ parva che.}
\end{quote}

\begin{quote}
Diwali is the festival of light.
\end{quote}

A neuter noun.

\begin{cousinweb}
\textbf{\gu{પર્વ}} (\emph{parva}) is a tatsama, Sanskrit \emph{parvan}, \textquotedblleft{}a knot, a joint\textquotedblright{}, and then a fixed point in the year, the days of the moon's changes and a festival (Monier-Williams). It is neuter: \textbf{\gu{પ્રકાશનું} \gu{પર્વ}}.
\end{cousinweb}

\begin{grammarlens}[title={What you've built}]
One more word for telling what happened.
\end{grammarlens}

\subsection*{Your turn}
\begin{itemize}
\raggedright
  \item \emph{Say it:} \emph{parva}
  \item \emph{Say it:} \emph{parva}, once more
  \item \emph{Say it:} \emph{divāḷī prakāśnũ parva che.} — Diwali is the festival of light.
  \item \emph{Recall:} say the Gujarati for a tradition, then the Gujarati for a custom, then say \emph{parva} again
\end{itemize}

\begin{tcolorbox}[breakable,colback=teal!4,colframe=teal!35!black,title={Before you move on}]
What does \gu{પર્વ} mean? (\textquotedblleft{}A holy day; a festival day\textquotedblright{}.) Say it once more.
\end{tcolorbox}
